\section*{General Project Guidelines}

\begin{tcolorbox}[colback=gray!6,colframe=gray!50!black]
The following guidelines apply to all projects.
\end{tcolorbox}

\subsection*{Scope \& Spirit}
Projects are meant to deepen your understanding of GPU computing through a concrete, measurable task.
You may \textbf{deviate} from the suggested items if you have a better idea, provided it stays within the \emph{spirit} of the project (same learning goals).
If you wish to propose a new scoring item not listed, \textbf{ask for approval} first; I will grant it if it has clear pedagogical value.

\subsection*{Deliverables}
\begin{enumerate}[label=\arabic*.]
    \item \textbf{Repository (required).} CUDA/C++ code with a clear structure and clear build/run instructions (e.g., in a \code{README.md}). Ideally, an automated build system (e.g., \code{CMakeLists.txt}).
    \item \textbf{Report (required).} Either \emph{Markdown} (in the repo) or a \emph{PDF} (3--10 pages). It should focus on:
    \begin{itemize}
        \item What you implemented (which items, with brief descriptions).
        \item Performance discussion: methodology, measurements, what limits throughput.
        \item Design/implementation choices: data layout, synchronization methods, trade-offs.
        \item Anything worth a discussion (e.g., failed attempts, unexpected findings).
    \end{itemize}
\end{enumerate}
Please note that the report is as important as the code for the assessment.

\subsection*{Submission \& Deadlines}
\begin{tcolorbox}[colback=white,colframe=black!10,title=Deadlines]
\textbf{Mid-semester feedback (optional):} \emph{July 17, 2026} (Sun), 23:59 CET.\\
\textbf{Hard deadline:} \emph{July 31, 2026} (Sun), 23:59 CET.
\end{tcolorbox}

\paragraph{Mid-semester feedback (optional).}
Submit a draft by July 17 to get written feedback from me:
\begin{itemize}
    \item Are you on track to reach the threshold (bonus) if submitted as is?
    \item Are you heading in a sensible technical direction?
    \item Any major flaw or mismatch with course goals?
    \item If your chosen direction is off-topic for the course, I will say so here.
\end{itemize}

\subsection*{Submission Workflow}
\begin{enumerate}[label=\arabic*.]
    \item Host the code on the LRZ GitLab. \textbf{Fork} the course repo (or create your own), open a \textbf{Merge Request}, and \textbf{assign me as reviewer}. If anything special should be noted you can add a short description in the MR.
    \item If you do not have LRZ GitLab access, email me a link to your repository.
\end{enumerate}

\subsection*{Evaluation}
\begin{itemize}
    \item \textbf{Grading.} Evaluation is \emph{pass/fail} based on the sum of points from the project's item list.
        Each item has a stated value; you can earn \emph{full}, \emph{partial}, or \emph{no} credit depending on the quality of your implementation and evidence.
        If your total is \textbf{$\geq$ 50 points}, you \textbf{pass} and receive a \textbf{--0.3} bonus on the final grade; otherwise, no bonus is awarded.
        The projects' maximum amount of points \textbf{exceeds} 50 and you do \textbf{not} need to implement everything and should rather pick the items that best fit your interests and time.
        Please mention in the report which items you implemented.
        Items completed with more depth or quality may receive \emph{extra} points at my discretion.
    \item \textbf{Workload expectation.} The bonus is calibrated for \(\approx\)20 hours of effective work \emph{per student}. Your submission should reflect this depth (code, measurements, and report).
    \item \textbf{What I value.} Code correctness first, then correct reasoning, then performance.
        Please note that it is still interesting if an optimization attempt \emph{fails} to deliver speedup, moreso if you can explain \emph{why}.
    \item \textbf{Evidence.} Include timing tables, plots, or other measurements for all relevant items.
        Detail as much as possible your measurement methodology (e.g., warmup, averaging, input sizes).
        Performance-related items will not be awarded points without evidence.
\end{itemize}

\subsection*{General Rules}
\begin{itemize}
    \item Your work will be assessed individually. You may collaborate with classmates for discussion, but the code and report must be your own.
    \item The use of AI tools is permitted but you will be held accountable for your submission.
    \item Unless specified otherwise, you are expected to propose pure CUDA kernels without using external libraries (e.g., Thrust, cuBLAS).
            Occasional, well-justified library calls are acceptable; in particular, in-kernel CUB collectives (e.g., \code{cub::WarpScan}, \code{cub::BlockScan}, \code{cub::BlockReduce}) are permitted.
            Avoid relying on host-side opaque primitives (e.g., \code{thrust::exclusive\_scan}, \code{cub::DeviceScan}) for your solution, although it is recommended to compare against them when appropriate (and it is perfectly fine to be slower).
            The objective of this project is to demonstrate kernel-level design (memory hierarchy, synchronization, divergence).
\end{itemize}

\begin{tcolorbox}[colback=gray!6,colframe=gray!50!black,title=Quick Checklist]
\begin{itemize}
    \item The \emph{readme} or \emph{report} explains how to compile and run your code.
    \item Following the instructions indeed builds and runs the code.
    \item There is a report (Markdown or PDF) covering implementation, measurements, and analysis.
    \item Submission via GitLab MR (assigned to me) or, if unavailable, by email with a link.
\end{itemize}
\end{tcolorbox}
